\section{Core Concepts in Metrology}

As previously mentioned, metrology is the science of measurement. As will become clear later, a measurement cannot exist without precision, and metrology is therefore closely connected to \textbf{uncertainty}.

To explore the world of metrology, several core concepts must first be defined and described:

\textbf{True Value:}  
This is the value that a physical object would have if it were observed completely correctly, without any measurement error. We may attempt to map or characterize the parameters of an object, but we can never obtain a perfect description. Therefore, the true value is an ideal concept. It exists in reality, since an object has the form and properties that it has, but due to complexity and limitations in our description of the object, we cannot determine the true value exactly.

\textbf{Measurement / Calibration:}  
A procedure and one or more actions that attempt to assess an instrument’s or object’s relationship to one or more true values of a physical quantity, as well as the potential deviation from that true value.

\textbf{Uncertainty:}  
Uncertainty can refer to many things, but as the word itself suggests, in this context it describes the level of confidence associated with a value and its connection to a physical quantity - or rather, the lack of complete certainty we have. In practice, this means a measured value together with a range within which the true value is believed to lie.

\textbf{Standard:}  
A standard (sometimes called a \textit{measurement standard}) is an object used to compare other quantities against. The standard is known with high precision, and therefore its value is trusted to a high degree. It is worth noting that a standard can also be an instrument. In fact, when you take your favorite voltmeter and measure a battery, the voltage measurement you obtain is referenced to the internal voltage reference of the voltmeter.

And thus, from this it should be clear that when we aim to collect data for IEEE, we are in fact probing reality using instruments. There are several best practices to follow to ensure proper measurements are taken. At the core of these are:

\begin{enumerate}
    \item Measurements cannot stand alone; an uncertainty must always accompany the reported value, together with the associated confidence level.
    \item Measurements must be documented and described to such a degree that any other person within your field should be capable of recreating your results, within the bounds of the stated uncertainty - always.
\end{enumerate}

\subsection{Responsibility of the Experimenter}

Up until this point, most students at university have not been given direct responsibility for measurements that may later be used by others. In this context, however, the situation is different.

Metrology relies heavily on a trickle-down effect. Data and measurements are published, reused, and built upon by other researchers and engineers. The measurements you perform today may form the foundation for future research, engineering designs, or automated systems.

In the context of IEEE data collection, the datasets produced may be used to train artificial intelligence systems that can have significant real-world impact. As the person responsible for performing the measurements, you must therefore ensure that the reported measurement values and their associated uncertainties are correct and valid.

If data is published with unrealistically small uncertainties or exaggerated accuracy, systems using that data
- such as AI models - may assign it far more confidence than it deserves. In applications such as autonomous vehicles, drones, or safety-critical systems, this can lead to incorrect decisions with real-world consequences.

Accurate documentation and honest uncertainty estimation are therefore not merely academic requirements; they are part of the professional and ethical responsibility of every engineer performing measurements. Should you later discover that incorrect data has been published, it must be corrected and those using the data must be informed. 

It is important to note that this is not simply a best practice or a suggestion for students - it is an ethical obligation that every engineer carries throughout their professional work.